\section{强制措施}
\begin{enumerate}
  \item \bs{仅限侦查阶段采取}！！公检法机关 \yjt 监委没有！！；针对嫌疑人/被告人；人身自由（\al{查扣冻 \cuo}）；预防性，临时性。
        \al{审判阶段，检可以对甲适用取保候审 \cuo}。
  \item 适用原则。
        \begin{enumerate}
          \item 法定性原则。\zd{法律事先作出明确的规定}。
          \item 必要性原则。不到万不得已不要用。能不罚就不罚，能少罚就少罚。
          \item 相当性/比例性原则。
                \begin{enumerate}
                  \item 对公民权利侵害最小。
                  \item \zd{严厉程度与危险性成正比}。\al{强奸、非法拘禁、判 10 年以上等 \dui}。
                \end{enumerate}
          \item 变更性原则。\zd{必要性减弱 \hz 消除}。
                \begin{enumerate}
                  \item \al{不能在法律规定的期限内办结的，予以释放 \dui{释放是变更，变更为其他也是变更}}。
                  \item \al{不能在法律规定的期限内办结的，检应当释放 \cuo}。
                  \item 有关才有放！\al{甲被取保候审，后应释放甲 \cuo}。
                \end{enumerate}
        \end{enumerate}
  \item 根据犯罪的轻重、证据的掌握情况、个人情况具体适用。
  \item \zx{vs 扭送}。\zd{普通公民针对嫌疑人/被告人强行带至公安司法机关的手段}。
        \begin{enumerate}
          \item 正在实行犯罪 \hz 犯罪后立即被发觉。
          \item 通缉在案、越狱逃跑、正在被追捕。
          \item 普通公民扭送一定要严加限制。
        \end{enumerate}
\end{enumerate}
